% ============================================================
% Chapter VI — Conclusion
% ============================================================
\section{Conclusion}

\subsection{Summary of Results}

This study predicted Formula One race finishing tiers (Podium, Point Scorer,
No Points) using within-race telemetry data from the 2023 Bahrain, Saudi
Arabian, and Australian Grands Prix. A dataset of 2,880 lap-level
observations was constructed from FastF1 telemetry, with features engineered
to represent the physical dynamics of race strategy: tire age, a linear fuel
burn-off proxy, their mean-centered interaction term, and tire compound
dummies, with grid position as a baseline control.

Box's M test rejected the homogeneity of covariance matrices assumption
($p < 2.2 \times 10^{-16}$), ruling out discriminant analysis and
motivating multinomial logistic regression as the primary inferential model.
VIF diagnostics confirmed the absence of problematic multicollinearity (all
VIF $< 3.5$). The MLE model identified GridPosition ($p = 0.000$) and the
Degradation $\times$ Fuel interaction term ($p = 0.0095$ for Podium) as
statistically significant predictors. Leave-one-group-out cross-validation
across the three races yielded 62.36\% overall accuracy, with podium recall
reaching 89\%.

\subsection{Connection to Literature}

The findings are consistent with and extend several threads in the existing
literature. The dominance of GridPosition as a predictor aligns with prior
F1 prediction studies \parencite{kuhn2022f1, song2020} that identified
starting position as the primary determinant of finishing order. The
rejection of covariance homogeneity mirrors findings by
\textcite{mendoza2023} in educational performance classification, suggesting
this pattern is general across multi-tier classification problems.

The central contribution---the significance of the Degradation $\times$ Fuel
interaction term---extends the tire degradation literature
\parencite{perdomo2021tire, baldi2022} and fuel load research
\parencite{oldknow2022fuel} by demonstrating that the \textit{coupling} of
these two forces, not just their individual effects, carries statistically
measurable predictive power. This supports the feature engineering philosophy
of \textcite{santos2020} and provides empirical backing for the physical
interaction models used in race strategy simulations.

The LOGO-CV validation approach responds directly to the methodological
concerns raised by \textcite{varoquaux2019} about overly optimistic
performance estimates in grouped data. By achieving 62.36\% accuracy under
strict cross-circuit validation, the study demonstrates that the statistical
relationships are genuine and not artifacts of data leakage.

\subsection{Future Research Directions}

\begin{enumerate}
  \item \textbf{Expand the race calendar.} Including all 23 races from a full
    F1 season would provide greater variance in weather conditions, circuit
    types, and team performance trajectories, strengthening both statistical
    power and generalizability.

  \item \textbf{Incorporate weather variables.} Rain probability, track
    temperature, and ambient temperature directly affect tire degradation
    rates and compound selection. These would capture a significant source of
    within-race variance.

  \item \textbf{Model non-linear tire degradation.} The linear tire age
    predictor could be replaced with a piecewise or polynomial term to
    capture the ``cliff'' effect where tire performance drops sharply after
    a critical wear threshold \parencite{baldi2022}.

  \item \textbf{Include traffic and DRS effects.} Lap time delta to the car
    ahead, DRS activation history, and pit stop timing would help explain
    mid-field variance and improve Point Scorer tier predictions.

  \item \textbf{Benchmark alternative classifiers.} Random forests,
    gradient-boosted trees, and ordinal logistic regression should be
    compared against the multinomial baseline to determine whether more
    flexible or order-aware models improve classification.

  \item \textbf{Apply to other racing series.} The feature engineering
    framework---fuel proxy, tire degradation, interaction terms, grid
    position control---is transferable to Formula 2, IndyCar, WEC, and
    other series with telemetry data.
\end{enumerate}

\subsection{Recommendations}

\begin{enumerate}
  \item F1 analytics teams should model tire degradation and fuel mass as an
    interaction rather than as independent effects when building race strategy
    simulations and real-time prediction systems.

  \item Researchers working with grouped motorsport data should adopt
    LOGO-CV rather than random $k$-fold cross-validation to produce honest
    out-of-sample performance estimates.

  \item Feature engineering in sports analytics should prioritize physically
    meaningful interaction terms over raw variables, as demonstrated by the
    Degradation $\times$ Fuel term's significance even after controlling for
    baseline pace.

  \item Future studies should incorporate dynamic race state variables
    (safety car status, DRS availability, tire compound age distribution
    across the field) to improve mid-field prediction accuracy.
\end{enumerate}

The null hypothesis---that tire degradation, fuel mass, their interaction,
and tire compound selection have no significant effect on finishing tier
log-odds after controlling for grid position---is rejected. The Degradation
$\times$ Fuel interaction term ($p = 0.0095 < \alpha = 0.05$) provides
statistically significant evidence that the physical coupling of tire wear
and fuel burn-off is a distinct, measurable predictor of podium finishes in
Formula One.
